\begin{table*}[!t]
\centering
\caption{Representative false-positive edges from the held-out audit. Text is shortened only at sentence boundaries. The examples illustrate three distinct extraction failure modes.}
\label{tab:edge_cases}
\small
\setlength{\tabcolsep}{5pt}
\renewcommand{\arraystretch}{1.12}
\begin{tabularx}{\textwidth}{@{}p{1.85cm} >{\raggedright\arraybackslash}X >{\raggedright\arraybackslash}X p{3.15cm}@{}}
\toprule
\textbf{Edge type} & \textbf{Earlier step} & \textbf{Later step} & \textbf{Why the edge is spurious} \\
\midrule
Variable reference & ``For $x=1$, plug in each integer root to determine $a$.'' & ``Thus, the possible values of $a$ are $\{-71,-27,-11,9\}$.'' & A shared symbol is insufficient to identify which case supplies the conclusion. \\
\addlinespace[2pt]
Order fallback & ``The quarters total $10\times\$0.25=\$2.50$.'' & ``The dimes total $12\times\$0.10=\$1.20$.'' & Adjacent calculations are parallel branches, not a support relation. \\
\addlinespace[2pt]
Expression reference & ``One hundred patches cost $100\times\$1.25=\$125$.'' & ``Selling 100 patches yields $100\times\$12=\$1{,}200$.'' & Reusing the count 100 does not make the cost calculation a premise of revenue. \\
\bottomrule
\end{tabularx}
\end{table*}
